\documentclass[11pt,a4paper]{article}
% Preambulo comun de los reportes de Tecnologias Celulares.
% Cada main.tex hace: \documentclass ... % Preambulo comun de los reportes de Tecnologias Celulares.
% Cada main.tex hace: \documentclass ... % Preambulo comun de los reportes de Tecnologias Celulares.
% Cada main.tex hace: \documentclass ... \input{../comun/preambulo.tex}
\usepackage[utf8]{inputenc}
\usepackage[T1]{fontenc}
\usepackage[spanish,mexico]{babel}
\usepackage{graphicx}
\usepackage[margin=2.5cm]{geometry}
\usepackage{amsmath}
\usepackage{amssymb}
\usepackage{booktabs}
\usepackage{tabularx}
\usepackage{array}
\usepackage{float}
\usepackage{setspace}
\usepackage{parskip}
\usepackage{fancyhdr}
\usepackage{textcomp}
\usepackage{xcolor}
\usepackage{caption}
\usepackage{fvextra}
\usepackage[hidelinks]{hyperref}

\onehalfspacing
\setlength{\parskip}{6pt}
\sloppy
\setlength{\emergencystretch}{3em}
\setlength{\headheight}{40pt}
\captionsetup{font=small,labelfont=bf}
\setlength{\tabcolsep}{6pt}

\pagestyle{fancy}
\fancyhf{}
\fancyhead[L]{\includegraphics[width=1cm]{/home/yovick/repos/UAZ-IEI/Assets/UIE.png}}
\fancyhead[R]{\small Tecnologías Celulares}
\fancyfoot[C]{\thepage}
\renewcommand{\headrulewidth}{0.4pt}

\DefineVerbatimEnvironment{codigo}{Verbatim}{%
  breaklines=true,breakanywhere=true,fontsize=\small,frame=single,
  framesep=3pt,rulecolor=\color{gray!50}}

% Figura centrada con pie: \figura[ancho]{archivo}{leyenda}
\newcommand{\figura}[3][0.85]{%
  \begin{center}
    \includegraphics[width=#1\linewidth]{#2}\\[2pt]
    {\small #3}
  \end{center}}

% Portada: \portada{numero}{titulo de la practica}
\newcommand{\portada}[2]{%
\begin{titlepage}
  \thispagestyle{empty}
  \begin{center}
    \vspace*{1cm}
    \includegraphics[width=4.5cm]{/home/yovick/repos/UAZ-IEI/Assets/uaz.png}
    \vspace{0.8cm}

    {\large Universidad Autónoma de Zacatecas}\\[4pt]
    {\large Unidad Académica de Ingeniería Eléctrica}\\[4pt]
    {\normalsize Ingeniería Electrónica Industrial con especialidad en Telecomunicaciones}

    \vspace{1.8cm}
    {\LARGE\bfseries Práctica #1}\\[10pt]
    {\Large\bfseries #2}
  \end{center}

  \vspace{2.2cm}
  \begin{center}
    \begin{tabular}{rl}
      \textbf{Materia:} & Tecnologías Celulares\\
      \textbf{Docente:} & Dr. Víktor Iván Rodríguez Abdalá\\
      \textbf{Alumno:}  & Arnold Yovick Ríos Zaldívar\\
      \textbf{Fecha:}   & Noviembre de 2026\\
    \end{tabular}
  \end{center}
  \vfill
\end{titlepage}}

\usepackage[utf8]{inputenc}
\usepackage[T1]{fontenc}
\usepackage[spanish,mexico]{babel}
\usepackage{graphicx}
\usepackage[margin=2.5cm]{geometry}
\usepackage{amsmath}
\usepackage{amssymb}
\usepackage{booktabs}
\usepackage{tabularx}
\usepackage{array}
\usepackage{float}
\usepackage{setspace}
\usepackage{parskip}
\usepackage{fancyhdr}
\usepackage{textcomp}
\usepackage{xcolor}
\usepackage{caption}
\usepackage{fvextra}
\usepackage[hidelinks]{hyperref}

\onehalfspacing
\setlength{\parskip}{6pt}
\sloppy
\setlength{\emergencystretch}{3em}
\setlength{\headheight}{40pt}
\captionsetup{font=small,labelfont=bf}
\setlength{\tabcolsep}{6pt}

\pagestyle{fancy}
\fancyhf{}
\fancyhead[L]{\includegraphics[width=1cm]{/home/yovick/repos/UAZ-IEI/Assets/UIE.png}}
\fancyhead[R]{\small Tecnologías Celulares}
\fancyfoot[C]{\thepage}
\renewcommand{\headrulewidth}{0.4pt}

\DefineVerbatimEnvironment{codigo}{Verbatim}{%
  breaklines=true,breakanywhere=true,fontsize=\small,frame=single,
  framesep=3pt,rulecolor=\color{gray!50}}

% Figura centrada con pie: \figura[ancho]{archivo}{leyenda}
\newcommand{\figura}[3][0.85]{%
  \begin{center}
    \includegraphics[width=#1\linewidth]{#2}\\[2pt]
    {\small #3}
  \end{center}}

% Portada: \portada{numero}{titulo de la practica}
\newcommand{\portada}[2]{%
\begin{titlepage}
  \thispagestyle{empty}
  \begin{center}
    \vspace*{1cm}
    \includegraphics[width=4.5cm]{/home/yovick/repos/UAZ-IEI/Assets/uaz.png}
    \vspace{0.8cm}

    {\large Universidad Autónoma de Zacatecas}\\[4pt]
    {\large Unidad Académica de Ingeniería Eléctrica}\\[4pt]
    {\normalsize Ingeniería Electrónica Industrial con especialidad en Telecomunicaciones}

    \vspace{1.8cm}
    {\LARGE\bfseries Práctica #1}\\[10pt]
    {\Large\bfseries #2}
  \end{center}

  \vspace{2.2cm}
  \begin{center}
    \begin{tabular}{rl}
      \textbf{Materia:} & Tecnologías Celulares\\
      \textbf{Docente:} & Dr. Víktor Iván Rodríguez Abdalá\\
      \textbf{Alumno:}  & Arnold Yovick Ríos Zaldívar\\
      \textbf{Fecha:}   & Noviembre de 2026\\
    \end{tabular}
  \end{center}
  \vfill
\end{titlepage}}

\usepackage[utf8]{inputenc}
\usepackage[T1]{fontenc}
\usepackage[spanish,mexico]{babel}
\usepackage{graphicx}
\usepackage[margin=2.5cm]{geometry}
\usepackage{amsmath}
\usepackage{amssymb}
\usepackage{booktabs}
\usepackage{tabularx}
\usepackage{array}
\usepackage{float}
\usepackage{setspace}
\usepackage{parskip}
\usepackage{fancyhdr}
\usepackage{textcomp}
\usepackage{xcolor}
\usepackage{caption}
\usepackage{fvextra}
\usepackage[hidelinks]{hyperref}

\onehalfspacing
\setlength{\parskip}{6pt}
\sloppy
\setlength{\emergencystretch}{3em}
\setlength{\headheight}{40pt}
\captionsetup{font=small,labelfont=bf}
\setlength{\tabcolsep}{6pt}

\pagestyle{fancy}
\fancyhf{}
\fancyhead[L]{\includegraphics[width=1cm]{/home/yovick/repos/UAZ-IEI/Assets/UIE.png}}
\fancyhead[R]{\small Tecnologías Celulares}
\fancyfoot[C]{\thepage}
\renewcommand{\headrulewidth}{0.4pt}

\DefineVerbatimEnvironment{codigo}{Verbatim}{%
  breaklines=true,breakanywhere=true,fontsize=\small,frame=single,
  framesep=3pt,rulecolor=\color{gray!50}}

% Figura centrada con pie: \figura[ancho]{archivo}{leyenda}
\newcommand{\figura}[3][0.85]{%
  \begin{center}
    \includegraphics[width=#1\linewidth]{#2}\\[2pt]
    {\small #3}
  \end{center}}

% Portada: \portada{numero}{titulo de la practica}
\newcommand{\portada}[2]{%
\begin{titlepage}
  \thispagestyle{empty}
  \begin{center}
    \vspace*{1cm}
    \includegraphics[width=4.5cm]{/home/yovick/repos/UAZ-IEI/Assets/uaz.png}
    \vspace{0.8cm}

    {\large Universidad Autónoma de Zacatecas}\\[4pt]
    {\large Unidad Académica de Ingeniería Eléctrica}\\[4pt]
    {\normalsize Ingeniería Electrónica Industrial con especialidad en Telecomunicaciones}

    \vspace{1.8cm}
    {\LARGE\bfseries Práctica #1}\\[10pt]
    {\Large\bfseries #2}
  \end{center}

  \vspace{2.2cm}
  \begin{center}
    \begin{tabular}{rl}
      \textbf{Materia:} & Tecnologías Celulares\\
      \textbf{Docente:} & Dr. Víktor Iván Rodríguez Abdalá\\
      \textbf{Alumno:}  & Arnold Yovick Ríos Zaldívar\\
      \textbf{Fecha:}   & Noviembre de 2026\\
    \end{tabular}
  \end{center}
  \vfill
\end{titlepage}}


\begin{document}

\portada{6}{Demodulación de señales OFDM}

\textbf{Script:} \texttt{practica6.py} (Python con NumPy/SciPy)

\section{Objetivo}

Implementar un receptor OFDM que recupere los símbolos transmitidos mediante la FFT,
demodule QPSK y evalúe el desempeño del sistema a través de la tasa de error de bit para
distintos valores de relación señal a ruido.

\section{Fundamento teórico}

El receptor OFDM aplica la FFT a cada símbolo recibido:
\begin{equation}
  Y[k] = \sum_{n=0}^{N-1} r[n]\,e^{-j2\pi kn/N} = X[k] + N[k],
\end{equation}
donde $N[k]$ es el ruido proyectado sobre la subportadora $k$. Gracias a la ortogonalidad,
cada subportadora se demodula por separado y la decisión QPSK se toma comparando las
partes real e imaginaria con cero.

La tasa de error de bit se estima por Monte Carlo:
\begin{equation}
  \mathrm{BER} = \frac{N_e}{N_b},
\end{equation}
con $N_e$ bits erróneos y $N_b$ bits transmitidos. Para QPSK con mapeo Gray, la BER
teórica en canal AWGN es
\begin{equation}
  P_b = Q\!\left(\sqrt{\frac{2E_b}{N_0}}\right).
\end{equation}
Si la energía por símbolo es $E_s=1$, entonces $E_b=E_s/2$ y la varianza total del ruido
complejo por muestra se fija en $N_0=E_b/(E_b/N_0)$.

\section{Desarrollo experimental}

\subsection{Parámetros}

\begin{tabularx}{\linewidth}{@{}*{2}{>{\raggedright\arraybackslash}X}@{}}
  \toprule
  Parámetro & Valor \\
  \midrule
  Bits por trama & 4096 \\
  Tramas por punto de SNR & 50 (204\,800 bits) \\
  Subportadoras & 64 \\
  Modulación & QPSK \\
  Canal & AWGN complejo \\
  $E_b/N_0$ evaluados & 0 a 14 dB en pasos de 2 dB \\
  \bottomrule
\end{tabularx}

\subsection{Procedimiento}

\begin{enumerate}
  \item Se generan 4096 bits aleatorios y se modulan en QPSK.
  \item Se aplica la IFFT por bloques de 64 y se serializa la señal.
  \item Se añade ruido gaussiano complejo con la varianza necesaria para el $E_b/N_0$
        deseado.
  \item En recepción se reordena la señal en bloques de 64, se aplica la FFT y se decide
        bit por bit con el signo de $I$ y $Q$.
  \item Se repite para 50 tramas por punto y se acumulan los errores.
\end{enumerate}

\section{Código}

\begin{codigo}
def recibir(rx, nfft=64):
    y = np.fft.fft(rx.reshape(-1, nfft).T, axis=0) / np.sqrt(nfft)
    return y.flatten(order="F")

def demodular_qpsk(simbolos):
    bits = np.empty(2 * len(simbolos), dtype=int)
    bits[0::2] = np.real(simbolos) > 0
    bits[1::2] = np.imag(simbolos) > 0
    return bits

sigma2 = 1 / (2 * 10 ** (ebn0_db / 10))   # Es = 1, Eb = Es/2
\end{codigo}

\section{Resultados}

\subsection{Constelaciones}

\figura{figs/constelacion.png}{Constelación transmitida y constelación recibida con $E_b/N_0=10$ dB.}

Con 10 dB la nube de ruido rodea cada punto de la constelación sin cruzar los ejes, por lo
que la mayoría de los símbolos se decide correctamente. El error vectorial normalizado
(RMS) medido fue de \textbf{0.224}, consistente con una relación $E_s/N_0=13$ dB.

\subsection{Tasa de error de bit}

\begin{tabular}{@{}ccc@{}}
  \toprule
  $E_b/N_0$ (dB) & BER simulada & BER teórica \\
  \midrule
  0  & 0.07912 & 0.07865 \\
  2  & 0.03772 & 0.03751 \\
  4  & 0.01289 & 0.01250 \\
  6  & 0.002412 & 0.002388 \\
  8  & 0.000200 & 0.000191 \\
  10 & $4.88\times10^{-6}$ & $3.87\times10^{-6}$ \\
  12 & 0 & $9.0\times10^{-9}$ \\
  14 & 0 & $6.8\times10^{-13}$ \\
  \bottomrule
\end{tabular}

\figura{figs/ber.png}{Curva de BER del sistema OFDM frente a la curva teórica de QPSK.}

La simulación sigue la curva teórica hasta 10 dB. A partir de 12 dB no se registran
errores en 204\,800 bits, lo que fija una cota superior de $4.9\times10^{-6}$, compatible
con los valores teóricos del orden de $10^{-8}$.

\section{Análisis de resultados}

\begin{itemize}
  \item \textbf{¿Cómo afecta el ruido a la constelación?} Convierte cada punto en una nube
        gaussiana centrada en el símbolo transmitido. Al aumentar el ruido, la nube crece y
        algunos símbolos cruzan los ejes, produciendo errores de bit.
  \item \textbf{¿Qué ocurre con la BER al incrementar la SNR?} Disminuye de forma
        aproximadamente exponencial: pasa de $7.9\times10^{-2}$ a 0 dB a
        $4.9\times10^{-6}$ a 10 dB. La pendiente coincide con la de QPSK en AWGN.
  \item \textbf{¿Por qué la FFT permite separar las subportadoras?} Porque las funciones
        base son ortogonales: la FFT proyecta la señal recibida sobre cada una de ellas y
        anula la contribución de las demás. Es la operación inversa de la IFFT del
        transmisor.
  \item \textbf{¿Qué diferencias existirían si las subportadoras no fueran ortogonales?}
        Aparecería interferencia entre subportadoras (ICI): la energía de una se sumaría
        como ruido en las vecinas, la constelación se dispersaría y la BER tendría un piso
        que no bajaría al aumentar la potencia.
\end{itemize}

\section{Cuestionario}

\begin{enumerate}
  \item \textbf{¿Cuál es la función de la FFT en un receptor OFDM?} Separar las
        subportadoras ortogonales y entregar una muestra compleja por subportadora, que es
        el símbolo transmitido más el ruido del canal.
  \item \textbf{¿Qué representa la BER?} La fracción de bits que el receptor decide
        incorrectamente respecto del total transmitido; es una medida directa de la
        confiabilidad del enlace.
  \item \textbf{¿Por qué OFDM utiliza múltiples subportadoras?} Para dividir un flujo de
        alta velocidad en muchos flujos lentos, alargar el símbolo y hacer que cada
        subcanal sea aproximadamente plano, lo que simplifica la ecualización y mejora la
        robustez frente a la dispersión.
  \item \textbf{¿Cómo se interpreta una constelación QPSK recibida?} Cada punto debería
        estar en una de las cuatro esquinas; la dispersión alrededor de cada esquina mide
        el ruido y la distorsión. Si la nube cruza los ejes, esos símbolos se deciden mal.
  \item \textbf{¿Qué factores limitan el desempeño de un sistema OFDM?} El ruido térmico,
        la interferencia entre subportadoras por desajustes de frecuencia o fase, la
        dispersión del canal más allá del prefijo cíclico, la no linealidad del amplificador
        (PAPR) y los errores de sincronización.
\end{enumerate}

\section{Conclusiones}

Se implementó un receptor OFDM completo y se midió su BER para $E_b/N_0$ de 0 a 14 dB. La
curva simulada coincidió con la teoría de QPSK en AWGN: $7.91\times10^{-2}$ a 0 dB,
$2.41\times10^{-3}$ a 6 dB y $4.88\times10^{-6}$ a 10 dB. La constelación recibida mostró
la nube de ruido esperada y la FFT recuperó correctamente las 64 subportadoras. El
resultado confirma que un enlace OFDM ideal, sin dispersión ni desajustes, se comporta
como un conjunto de canales QPSK paralelos.

\end{document}
